% 第九卷 15.1.2 节：卷积下采样前端与Conformer块（跨栏，约140 mm宽）。
% 蓝色：输入特征；白底黄描边：模型部件；浅黄底：卷积模块；红色：编码输出。
\begin{tikzpicture}[
  x=1mm,y=1mm,
  figure picture,
  inp/.style={figure node,draw=FigureInput,fill=FigureInputFill,
    minimum height=7mm,minimum width=46mm,inner xsep=2pt,inner ysep=2pt,
    font=\figurefont\scriptsize},
  mod/.style={figure node,draw=FigureModel,fill=FigureModelFill,
    minimum height=7mm,minimum width=46mm,inner xsep=2pt,inner ysep=2pt,
    font=\figurefont\scriptsize},
  hot/.style={mod,fill=FigureCream},
  sub/.style={figure node,draw=FigureModel,fill=FigureModelFill,
    minimum height=6.5mm,minimum width=40mm,inner xsep=2pt,inner ysep=2pt,
    font=\figurefont\scriptsize},
  res/.style={figure node,draw=FigureOutput,fill=FigureOutputFill,
    minimum height=7mm,minimum width=46mm,inner xsep=2pt,inner ysep=2pt,
    font=\figurefont\scriptsize},
  plus/.style={circle,draw=FigureStroke,fill=white,line width=.6pt,
    inner sep=0pt,minimum size=3.6mm,font=\figurefont\scriptsize},
  note/.style={font=\figurefont\scriptsize,text=FigureMuted,anchor=west},
  fl/.style={figure flow},
  skip/.style={figure flow,draw=FigureMuted,line width=.6pt},
  guide/.style={figure link,dash pattern=on 1.5pt off 1.5pt}
]
  % ---------- 前端 ----------
  \node[inp] (x) at (0,0) {log-Mel特征$\bm X$};
  \node[note] at (25,0) {$T_0\times80$};
  \node[mod] (cv) at (0,12) {两层$3\times3$卷积，步长2，ReLU};
  \node[note] at (25,12) {$c\times\lceil T_0/4\rceil\times20$};
  \node[mod] (ln0) at (0,24) {重排、线性投影、Dropout};
  \node[note] at (25,24) {$T\times d$};
  \draw[fl] (x) -- (cv);
  \draw[fl] (cv) -- (ln0);

  % ---------- Conformer块 ----------
  \draw[figure subgroup] (-33,31) rectangle (24,111);
  \node[font=\figurefont\scriptsize,text=FigureMuted,anchor=north west] at (-32.5,110.5) {$\times N$};

  \node[mod] (f1) at (0,39) {前馈模块（$d\to4d\to d$，Swish）};
  \node[plus] (p1) at (0,48) {$+$};
  \node[mod] (at) at (0,56) {多头自注意力（相对位置编码）};
  \node[plus] (p2) at (0,65) {$+$};
  \node[hot] (cm) at (0,73) {卷积模块};
  \node[plus] (p3) at (0,82) {$+$};
  \node[mod] (f2) at (0,90) {前馈模块};
  \node[plus] (p4) at (0,99) {$+$};
  \node[mod] (ln) at (0,106) {层归一化};
  \node[res] (y) at (0,120) {编码输出$\bm H$};
  \node[note] at (25,120) {$T\times d$，$T\approx T_0/4$};

  \draw[fl] (ln0) -- (f1);
  \draw[fl] (f1) -- node[right,font=\figurefont\scriptsize,text=FigureMuted]{$\times\frac12$} (p1);
  \draw[fl] (p1) -- (at);
  \draw[fl] (at) -- (p2);
  \draw[fl] (p2) -- (cm);
  \draw[fl] (cm) -- (p3);
  \draw[fl] (p3) -- (f2);
  \draw[fl] (f2) -- node[right,font=\figurefont\scriptsize,text=FigureMuted]{$\times\frac12$} (p4);
  \draw[fl] (p4) -- (ln);
  \draw[fl] (ln) -- (y);

  % 残差支路
  \foreach \a/\b in {33.5/p1,51.5/p2,68.5/p3,85.5/p4}{
    \draw[skip] (0,\a) -- (-28,\a) |- (\b.west);
  }

  % ---------- 卷积模块展开 ----------
  \draw[figure subgroup] (52,36) rectangle (128,112);
  \node[font=\figurefont\scriptsize,text=FigureMuted,anchor=north west] at (52.5,111.5) {卷积模块展开};
  \node[sub] (c1) at (75,42) {层归一化};
  \node[sub] (c2) at (75,52) {逐点卷积$d\to2d$};
  \node[sub] (c3) at (75,62) {GLU门控$2d\to d$};
  \node[sub] (c4) at (75,72) {深度卷积，核长$K$};
  \node[sub] (c5) at (75,82) {批归一化};
  \node[sub] (c6) at (75,92) {Swish};
  \node[sub] (c7) at (75,102) {逐点卷积$d\to d$，Dropout};
  \foreach \i/\j in {c1/c2,c2/c3,c3/c4,c4/c5,c5/c6,c6/c7}{\draw[fl] (\i) -- (\j);}
  \node[note] at (96,42) {$T\times d$};
  \node[note] at (96,52) {$T\times2d$};
  \node[note] at (96,62) {$T\times d$};
  \node[note] at (96,72) {逐通道沿时间滤波};
  \node[note] at (96,82) {批内统计量};
  \node[note] at (96,92) {$x\,\sigma(x)$};
  \node[note] at (96,102) {$T\times d$};

  \draw[guide] (cm.east) -- (52,36);
  \draw[guide] (cm.east) -- (52,112);
\end{tikzpicture}
